\documentclass[aspectratio=169,xcolor=dvipsnames]{beamer}
\usetheme{Berlin}

\usepackage[utf8]{inputenc}   % UTF-8 encoding for Spanish accents in text
\usepackage{hyperref}
\usepackage{graphicx}
\usepackage{booktabs}
\usepackage{amsmath}
\usepackage{lettrine}
\setbeamertemplate{caption}[numbered]
\usepackage[dvipsnames,svgnames,x11names]{xcolor}
\usepackage{xurl}
\usepackage{algorithm}
\usepackage{algorithmicx}
\usepackage{algpseudocode}
\usepackage{adjustbox}

\hypersetup{
    colorlinks=true,
    linkcolor=cyan,
    filecolor=blue,
    urlcolor=blue,
    citecolor=blue,
}

%----------------------------------------------------------------------------------------
%   CODE LISTINGS SETTINGS
%----------------------------------------------------------------------------------------
\usepackage{listings}
\definecolor{backcolour}{rgb}{0.97,0.97,0.99}
\definecolor{codegreen}{rgb}{0,0.6,0}

\lstdefinestyle{PythonStyle}{
  language=Python,
  basicstyle=\ttfamily\scriptsize,
  keywordstyle=\color{blue}\bfseries,
  commentstyle=\color{codegreen},
  stringstyle=\color{red},
  breaklines=true,
  numbers=left,
  numbersep=5pt,
  frame=lines,
  backgroundcolor=\color{backcolour},
  showstringspaces=false,
  morekeywords={True,False,None}
}
\lstset{style=PythonStyle}

%----------------------------------------------------------------------------------------
%   TITLE CONFIGURATION
%----------------------------------------------------------------------------------------
\title{Listas en Python}
\subtitle{Estructuras de datos unidimensionales (vectores)}
\author{Prof. D.Sc. BARSEKH-ONJI Aboud}
\institute{
    Facultad de Ingeniería \\
    Universidad Anáhuac México
}
\date{\today}

%----------------------------------------------------------------------------------------
%   AUTO AGENDA AT EACH SECTION
%----------------------------------------------------------------------------------------
\AtBeginSection[]{
  \begin{frame}{Agenda}
    \tableofcontents[currentsection]
  \end{frame}
}

\begin{document}

\begin{frame}
    \titlepage
\end{frame}

\begin{frame}{Agenda}
    \tableofcontents
\end{frame}

%------------------------------------------------
\section{¿Qué es una lista?}
%------------------------------------------------

\begin{frame}{¿Qué es una lista en Python?}
    \begin{block}{Definición}
        Una \textbf{lista} es una estructura de datos que almacena una colección ordenada de elementos. En Python, los elementos pueden ser de cualquier tipo: números, cadenas de texto (\textit{strings}), booleanos, etc.
    \end{block}
    \vspace{0.3cm}
    \begin{itemize}
        \item Es una estructura \textbf{mutable}: se puede modificar después de crearla.
        \item Los elementos están \textbf{indexados}: el primer elemento tiene índice \texttt{0}.
        \item Se puede almacenar elementos de \textbf{diferente tipo} en la misma lista.
        \item En este curso la usaremos como \textbf{vector unidimensional}.
    \end{itemize}
\end{frame}

\begin{frame}[fragile]{Crear una lista}
    \begin{block}{Sintaxis básica}
        \texttt{nombre\_lista = [elemento0, elemento1, elemento2, \ldots]}
    \end{block}
    \vspace{0.2cm}
\begin{lstlisting}[style=PythonStyle]
# Lista de enteros
notas = [85, 90, 72, 95, 60]
# Lista de strings (cadenas de texto)
nombres = ["Ana", "Luis", "Maria", "Carlos"]
# Lista de flotantes
precios = [10.5, 23.0, 8.75, 45.99]
# Lista vacia
vacia = []
print(notas)    # [85, 90, 72, 95, 60]
print(nombres)  # ['Ana', 'Luis', 'Maria', 'Carlos']
\end{lstlisting}
\end{frame}

\begin{frame}{Índices de una lista}
    \begin{block}{Índice positivo}
        El primer elemento tiene índice \textbf{0}, el segundo índice \textbf{1}, y así sucesivamente.
    \end{block}
    \vspace{0.3cm}
    \begin{center}
    \begin{tabular}{ccccc}
        \toprule
        \textbf{Índice} & 0 & 1 & 2 & 3 \\
        \midrule
        \textbf{Valor}  & ``Ana'' & ``Luis'' & ``Maria'' & ``Carlos'' \\
        \bottomrule
    \end{tabular}
    \end{center}
    \vspace{0.3cm}
    \begin{alertblock}{¡Cuidado!}
        Si intentas acceder a un índice fuera del rango, Python genera un error: \texttt{IndexError: list index out of range}.
    \end{alertblock}
\end{frame}

\begin{frame}[fragile]{Acceder a elementos de una lista}
\begin{lstlisting}[style=PythonStyle]
nombres = ["Ana", "Luis", "Maria", "Carlos"]
# Acceder por indice positivo
print(nombres[0])   # Ana
print(nombres[2])   # Maria
# Acceder por indice negativo
print(nombres[-1])  # Carlos (ultimo elemento)
print(nombres[-2])  # Maria  (penultimo elemento)
# Longitud de la lista
print(len(nombres)) # 4
\end{lstlisting}
    \begin{exampleblock}{Índices negativos}
        Python permite índices negativos: \texttt{-1} accede al último elemento, \texttt{-2} al penúltimo, etc.
    \end{exampleblock}
\end{frame}

%------------------------------------------------
\section{Operaciones básicas}
%------------------------------------------------

\begin{frame}[fragile]{Modificar elementos}
\begin{lstlisting}[style=PythonStyle]
notas = [85, 90, 72, 95, 60]
# Modificar un elemento por su indice
notas[2] = 80
print(notas)  # [85, 90, 80, 95, 60]
# Agregar un elemento al final
notas.append(88)
print(notas)  # [85, 90, 80, 95, 60, 88]
# Eliminar el ultimo elemento
notas.pop()
print(notas)  # [85, 90, 80, 95, 60]
# Eliminar un elemento por su indice
notas.pop(1)
print(notas)  # [85, 80, 95, 60]
\end{lstlisting}
\end{frame}

\begin{frame}[fragile]{Operaciones con \texttt{len}, \texttt{sum}, \texttt{min}, \texttt{max}}
\begin{lstlisting}[style=PythonStyle]
notas = [85, 90, 72, 95, 60]
# Numero de elementos
print(len(notas))   # 5
# Suma de todos los elementos
print(sum(notas))   # 402
# Valor minimo
print(min(notas))   # 60
# Valor maximo
print(max(notas))   # 95
# Promedio (suma / cantidad)
promedio = sum(notas) / len(notas)
print(promedio)     # 80.4
\end{lstlisting}
\end{frame}

\begin{frame}[fragile]{Recorrer una lista con \texttt{for}}
    \begin{block}{Patrón fundamental}
        Usamos un ciclo \texttt{for} para acceder a cada elemento de la lista uno por uno.
    \end{block}
\begin{lstlisting}[style=PythonStyle]
frutas = ["manzana", "naranja", "uva", "mango"]
# Iterar sobre los elementos directamente
for fruta in frutas:
    print(fruta)
# Iterar usando el indice
for i in range(len(frutas)):
    print(i, frutas[i])
# Salida:
# 0 manzana
# 1 naranja
# 2 uva
# 3 mango
\end{lstlisting}
\end{frame}

\begin{frame}[fragile]{Buscar un elemento: operador \texttt{in}}
\begin{lstlisting}[style=PythonStyle]
frutas = ["manzana", "naranja", "uva", "mango"]
# Verificar si un elemento esta en la lista
if "uva" in frutas:
    print("Si esta la uva")
else:
    print("No esta la uva")
# Encontrar el indice de un elemento
posicion = frutas.index("naranja")
print(posicion)   # 1
# Contar cuantas veces aparece un elemento
repeticiones = frutas.count("mango")
print(repeticiones)  # 1
\end{lstlisting}
\end{frame}

\begin{frame}[fragile]{Ordenar y otras operaciones útiles}
\begin{lstlisting}[style=PythonStyle]
numeros = [30, 10, 50, 20, 40]
# Ordenar de menor a mayor (modifica la lista original)
numeros.sort()
print(numeros)  # [10, 20, 30, 40, 50]
# Ordenar de mayor a menor
numeros.sort(reverse=True)
print(numeros)  # [50, 40, 30, 20, 10]
# Invertir el orden
numeros.reverse()
print(numeros)  # [10, 20, 30, 40, 50]
# Crear lista con range
pares = list(range(0, 11, 2))
print(pares)    # [0, 2, 4, 6, 8, 10]
\end{lstlisting}
\end{frame}

%------------------------------------------------
\section{Ejemplos resueltos}
%------------------------------------------------

\begin{frame}[fragile]{Ejemplo 1 -- Contar aprobados y reprobados}
    \begin{exampleblock}{Problema}
        Dada una lista de calificaciones, contar cuántos alumnos aprobaron (nota $\geq 70$) y cuántos reprobaron.
    \end{exampleblock}
\begin{lstlisting}[style=PythonStyle]
notas = [85, 60, 72, 45, 90, 68, 78, 55]
aprobados = 0
reprobados = 0
for nota in notas:
    if nota >= 70:
        aprobados = aprobados + 1
    else:
        reprobados = reprobados + 1
print("Aprobados: ", aprobados)   # Aprobados:  5
print("Reprobados:", reprobados)  # Reprobados: 3
\end{lstlisting}
\end{frame}

\begin{frame}[fragile]{Ejemplo 2 -- Calcular el promedio de una lista}
    \begin{exampleblock}{Problema}
        Calcular el promedio de una lista de números ingresados por el usuario.
    \end{exampleblock}
\begin{lstlisting}[style=PythonStyle]
# Crear lista vacia
datos = []
# Pedir 5 numeros al usuario
for i in range(5):
    valor = float(input("Ingresa un numero: "))
    datos.append(valor)
# Calcular el promedio
total = 0
for numero in datos:
    total = total + numero
promedio = total / len(datos)
print("Los datos:", datos)
print("Promedio:", promedio)
\end{lstlisting}
\end{frame}

\begin{frame}[fragile]{Ejemplo 3 -- Encontrar el mayor sin usar \texttt{max()}}
    \begin{exampleblock}{Problema}
        Encontrar el valor máximo de una lista \textbf{sin} usar la función \texttt{max()}.
    \end{exampleblock}
\begin{lstlisting}[style=PythonStyle]
numeros = [34, 12, 78, 5, 56, 90, 23]
# Suponemos que el mayor es el primer elemento
mayor = numeros[0]
for i in range(1, len(numeros)):
    if numeros[i] > mayor:
        mayor = numeros[i]
print("El valor maximo es:", mayor)  # 90
\end{lstlisting}
    \begin{block}{Idea clave}
        Se inicializa \texttt{mayor} con el primer elemento y se recorre el resto comparando uno a uno.
    \end{block}
\end{frame}

\begin{frame}[fragile]{Ejemplo 4 -- Filtrar elementos de una lista}
    \begin{exampleblock}{Problema}
        Dada una lista de nombres, crear una nueva lista con los nombres que empiezan con la letra ``A''.
    \end{exampleblock}
\begin{lstlisting}[style=PythonStyle]
nombres = ["Ana", "Bruno", "Alicia", "Carlos", "Alejandro", "Diana"]
# Lista para guardar los resultados
nombres_A = []
for nombre in nombres:
    if nombre[0] == "A":
        nombres_A.append(nombre)
print("Nombres con A:", nombres_A)
# Nombres con A: ['Ana', 'Alicia', 'Alejandro']
\end{lstlisting}
    \begin{alertblock}{Nota}
        \texttt{nombre[0]} accede al \textbf{primer carácter} del string, ya que un string también es indexable.
    \end{alertblock}
\end{frame}

\begin{frame}[fragile]{Ejemplo 5 -- Clasificar elementos en tres categorías}
    \begin{exampleblock}{Problema}
        Clasificar las calificaciones en: Excelente ($\geq 90$), Bueno ($70$--$89$) y Reprobado ($< 70$).
    \end{exampleblock}
\begin{lstlisting}[style=PythonStyle]
calificaciones = [95, 65, 88, 72, 55, 91, 78, 45, 83]
excelente = []
bueno = []
reprobado = []
for cal in calificaciones:
    if cal >= 90:
        excelente.append(cal)
    elif cal >= 70:
        bueno.append(cal)
    else:
        reprobado.append(cal)

print("Excelente:", excelente)   # [95, 91]
print("Bueno:",     bueno)       # [88, 72, 78, 83]
print("Reprobado:", reprobado)   # [65, 55, 45]
\end{lstlisting}
\end{frame}

%------------------------------------------------
\section{Ejercicios para resolver}
%------------------------------------------------

\begin{frame}{Ejercicios para resolver}
    \begin{alertblock}{Instrucciones}
        Resuelve cada ejercicio en Python usando listas y los ciclos \texttt{for} e instrucciones \texttt{if / elif / else} aprendidas. No uses funciones avanzadas fuera de las vistas en clase.
    \end{alertblock}
\end{frame}

\begin{frame}{Ejercicio 1 -- Suma de pares}
    \begin{block}{Enunciado}
        Dada la siguiente lista de enteros:
        \[ \texttt{numeros = [3, 8, 15, 22, 7, 14, 5, 18, 11, 6]} \]
        \begin{enumerate}
            \item Calcula la \textbf{suma} de todos los números pares de la lista.
            \item Imprime cuántos números pares hay en total.
        \end{enumerate}
    \end{block}
    \vspace{0.3cm}
    \textbf{Pista:} Un número es par si \texttt{numero \% 2 == 0}.
\end{frame}

\begin{frame}{Ejercicio 2 -- Lista invertida}
    \begin{block}{Enunciado}
        Dada la siguiente lista:
        \[ \texttt{original = [10, 20, 30, 40, 50]} \]
        \begin{enumerate}
            \item Crea una \textbf{nueva lista} llamada \texttt{invertida} que contenga los mismos elementos pero en orden inverso, \textbf{sin} usar \texttt{.reverse()} ni \texttt{[::-1]}.
            \item Imprime ambas listas.
        \end{enumerate}
    \end{block}
    \vspace{0.3cm}
    \textbf{Pista:} Recorre la lista original de atrás hacia adelante usando \texttt{range(len(original)-1, -1, -1)}.
\end{frame}

\begin{frame}{Ejercicio 3 -- Temperatura promedio}
    \begin{block}{Enunciado}
        Se tienen las temperaturas (en °C) registradas durante una semana:
        \[ \texttt{temps = [22.5, 25.0, 19.8, 30.1, 27.3, 21.0, 23.6]} \]
        \begin{enumerate}
            \item Calcula la \textbf{temperatura promedio} de la semana.
            \item Imprime los días (índice) donde la temperatura fue \textbf{mayor al promedio}.
            \item Indica cuántos días estuvieron \textbf{por debajo} del promedio.
        \end{enumerate}
    \end{block}
\end{frame}

\begin{frame}{Ejercicio 4 -- Contar vocales en una lista de nombres}
    \begin{block}{Enunciado}
        Dada la lista:
        \[ \texttt{nombres = ["Eduardo", "Iris", "Omar", "Ursula", "Pablo"]} \]
        \begin{enumerate}
            \item Para cada nombre, cuenta cuántas \textbf{vocales} (a, e, i, o, u) contiene (sin importar mayúsculas/minúsculas).
            \item Imprime el nombre junto con su cantidad de vocales.
        \end{enumerate}
    \end{block}
    \vspace{0.2cm}
    \textbf{Pista:} Puedes usar \texttt{nombre.lower()} para convertir a minúsculas y luego recorrer cada carácter con \texttt{for letra in nombre}.
\end{frame}

\begin{frame}{Ejercicio 5 -- Reemplazar negativos}
    \begin{block}{Enunciado}
        Dada la lista de mediciones:
        \[ \texttt{datos = [4.5, -2.1, 8.3, -0.5, 6.0, -3.7, 1.2]} \]
        \begin{enumerate}
            \item Recorre la lista y \textbf{reemplaza} cada valor negativo por \texttt{0.0}.
            \item Imprime la lista original (antes de modificar) y la lista resultante.
        \end{enumerate}
    \end{block}
    \vspace{0.2cm}
    \textbf{Pista:} Necesitas acceder por índice: \texttt{for i in range(len(datos))}.\\
    Usa \texttt{datos\_copia = datos.copy()} para guardar la original antes de modificar.
\end{frame}

\begin{frame}[fragile]{Resumen de operaciones vistas}
    \begin{columns}
        \column{0.5\textwidth}
        \textbf{Crear y modificar:}
        \begin{itemize}
            \item \texttt{lista = []}
            \item \texttt{lista.append(x)}
            \item \texttt{lista.pop(i)}
            \item \texttt{lista[i] = valor}
        \end{itemize}
        \vspace{0.2cm}
        \textbf{Información:}
        \begin{itemize}
            \item \texttt{len(lista)}
            \item \texttt{sum(lista)}
            \item \texttt{min(lista)} / \texttt{max(lista)}
            \item \texttt{lista.count(x)}
            \item \texttt{lista.index(x)}
        \end{itemize}

        \column{0.5\textwidth}
        \textbf{Recorrer:}
        \begin{itemize}
            \item \texttt{for x in lista:}
            \item \texttt{for i in range(len(lista)):}
        \end{itemize}
        \vspace{0.2cm}
        \textbf{Buscar / Ordenar:}
        \begin{itemize}
            \item \texttt{if x in lista:}
            \item \texttt{lista.sort()}
            \item \texttt{lista.sort(reverse=True)}
            \item \texttt{lista.reverse()}
            \item \texttt{lista.copy()}
        \end{itemize}
    \end{columns}
\end{frame}

\begin{frame}{}
    \centering
    \large Facultad de Ingeniería\\
    Universidad Anáhuac México\\[0.3cm]
    \normalsize \texttt{aboud.barsekh@anahuac.mx}
\end{frame}

\end{document}
